\section{Discusión de resultados}
\label{sec:discusion}

El control en \(q=1\) muestra que la transferencia, la construcción modal de
la semilla y la continuación recuperan el procedimiento entero de
Leonov--Kuznetsov \cite{LeonovKuznetsov2013,KuznetsovEtAl2017}. Para
\(0<q<1\), la adaptación conserva la misma idea de localización, pero distingue
dos problemas que cumplen funciones diferentes: Liouville--Weyl proporciona el
predictor armónico estacionario y el problema de valor inicial de Caputo
determina la trayectoria causal. Esta separación permite usar la información
frecuencial sin sustituir la memoria que forma parte del sistema fraccionario.

Las reconstrucciones bibliográficas establecen controles claros para comparar
las rutas propuestas. Con los parámetros nominales de Danca
\cite{Danca2017}, la integración ABM de \cref{eq:chua_no_suave} produce una
trayectoria postransitoria de apariencia periódica. La ruta propuesta para el
mismo sistema físico construye su propia semilla sesgada, la transporta con
memoria de Caputo y también alcanza un conjunto atrayente de apariencia
periódica, con una frecuencia dominante próxima a la del control
bibliográfico. Los diagnósticos de ambos cálculos indican un régimen regular. La
comparación adicional con
\((m_0,m_1)=(-0.2,-1.2)\) produce contactos desde vecindades de los
equilibrios externos y confirma que el clasificador distingue un conjunto
autoexcitado del caso localizado.

En la familia con arctangente, las condiciones iniciales y la recurrencia ADM
publicadas por Wu et al. \cite{Wu2023ArctanChua} producen dos trayectorias
simétricas de apariencia periódica. El caso \texttt{c590} de
\cref{eq:chua_arctan,eq:parametros_c590} procede de una exploración acotada en
\(q=1\), una selección cuantitativa y un refinamiento mediante problemas de
Caputo independientes. Su trayectoria postransitoria es aperiódica bajo el
filtro empleado. La comparación muestra que el régimen observado resulta de
la combinación entre la forma arctangente, los parámetros, la semilla y el
protocolo fraccionario.

\ifHAFOBasinMaterial
Los sondeos separan el alcance local predeclarado de la exploración
macroscópica. Para el sistema no suave, el protocolo local \(r\leq0.01\)
produjo \(0/7\,200\) contactos y cero fallos numéricos; éste es el soporte
finito para la compatibilidad local con ocultedad del conjunto regular. La
extensión macroscópica añadió 10\,200 sondas en \(r=0.03,0.1,0.3\):
\(0/6\,600\) contactos en los dos primeros radios y \(37/3\,600\) en el
tercero. El acumulado combinado es \(0/13\,800\) hasta \(r=0.1\); la afirmación
de ocultedad conserva el alcance local predeclarado. Para
\texttt{c590} se obtuvieron \(0/8\,400\) contactos hasta \(r=0.3\),
\(22/2\,400\) en \(r=1\) y \(588/2\,700\) en \(r=2\), con 131 divergencias
en este último bloque. El primer contacto aparece en \(r=1\); \(r=2\) es una
evaluación macroscópica adicional y no amplía el dominio local sin contactos.
\fi

La función descriptiva sesgada determina la semilla no suave mediante el
cierre armónico, la construcción modal y la continuación causal. En el caso
con arctangente, la exploración paramétrica y el refinamiento fraccionario
seleccionan la condición inicial. La evaluación posterior de cuenca emplea
radios, muestras y tolerancias previamente fijados; su alcance depende de
estos parámetros y del horizonte de integración.

\section{Conclusiones}
\label{sec:conclusiones}
\enlargethispage{2\baselineskip}

Se desarrollaron dos rutas teórico--numéricas para localizar conjuntos
atrayentes en sistemas de Chua fraccionarios. En
\cref{eq:chua_no_suave}, una función descriptiva sesgada construye la semilla
y una continuación causal de 21 etapas conserva una sola historia de Caputo.
En \cref{eq:chua_arctan}, una exploración acotada en orden entero y el
refinamiento mediante problemas de Caputo determinan el caso
\texttt{c590}. La formulación contiene como límite el procedimiento de
Leonov--Kuznetsov cuando \(q=1\).

Las reconstrucciones de Danca y Wu proporcionan controles bibliográficos
regulares. Las rutas propuestas producen, respectivamente, una trayectoria
postransitoria de apariencia periódica y otra aperiódica. En el caso no
suave, las perturbaciones de la condición inicial a paso fijo confirman
contacto con la nube objetivo; queda pendiente el refinamiento temporal
especificado en \cref{sec:control_numerico}.
\ifHAFOBasinMaterial
Los sondeos delimitan regiones muestreadas sin contacto y detectan acceso a
las nubes candidatas al aumentar el radio. La compatibilidad con ocultedad
es local y depende del protocolo finito; la separación de cuencas en un
continuo de condiciones iniciales requiere un argumento adicional.
\fi

Los sistemas enteros no Chua amplían la evaluación a Kalman--Fitts, MAVPD y
PLL. El mapa de conmutación, la continuación paramétrica y el retorno
cilíndrico permiten localizar conjuntos en modelos con obstrucciones
diferentes al balance armónico directo. La cobertura de sistemas
fraccionarios no Chua es todavía parcial y requiere los controles de
convergencia y cuenca indicados en \cref{sec:cobertura-cuatro-familias}.

Los PP exactos de PLL y MAVPD proporcionan semillas geométricas; los criterios
de ausencia de PP en modelos de tipo jerk delimitan las hipótesis bajo las
cuales ese recurso está disponible. Su condición algebraica no determina el
destino de la trayectoria. La utilidad topológica reside en complementar la
localización con un argumento sobre las cuencas: el ejemplo de
\cref{subsec:guide-hidden-cycle-topology} separa la existencia del ciclo de la
exclusión de vecindades de equilibrio. Esta distinción establece las
condiciones necesarias para convertir un candidato localizado en un
resultado de ocultedad.
